% !TEX program = xelatex
% Overleaf compiler: XeLaTeX
\documentclass[14pt,a4paper]{extarticle}

\usepackage{fontspec}
\usepackage{polyglossia}
\setdefaultlanguage{russian}
\setotherlanguage{english}
\IfFontExistsTF{Times New Roman}{%
    \setmainfont{Times New Roman}
    \newfontfamily\cyrillicfont{Times New Roman}[Script=Cyrillic]
}{%
    \IfFontExistsTF{Times New Roman Cyr}{%
        \setmainfont{Times New Roman Cyr}
        \newfontfamily\cyrillicfont{Times New Roman Cyr}[Script=Cyrillic]
    }{%
        \IfFontExistsTF{Tinos}{%
            \setmainfont{Tinos}
            \newfontfamily\cyrillicfont{Tinos}[Script=Cyrillic]
        }{%
            \IfFontExistsTF{Liberation Serif}{%
                \setmainfont{Liberation Serif}
                \newfontfamily\cyrillicfont{Liberation Serif}[Script=Cyrillic]
            }{%
                \setmainfont{FreeSerif}
                \newfontfamily\cyrillicfont{FreeSerif}[Script=Cyrillic]
            }%
        }%
    }%
}
\setmonofont{DejaVu Sans Mono}

\usepackage[a4paper,left=30mm,right=15mm,top=20mm,bottom=20mm]{geometry}
\usepackage{setspace}
\usepackage{indentfirst}
\usepackage{amsmath,amssymb}
\usepackage{graphicx}
\usepackage{float}
\usepackage{array}
\usepackage{longtable}
\usepackage{ragged2e}
\usepackage{booktabs}
\usepackage{caption}
\usepackage{hyperref}
\usepackage{xurl}
\usepackage{enumitem}
\usepackage{makecell}
\usepackage{multirow}
\usepackage{titlesec}
\usepackage{tocloft}
\usepackage[most]{tcolorbox}
\usepackage{fvextra}

\hypersetup{hidelinks}
\urlstyle{same}
\onehalfspacing
\setlength{\parindent}{1.25cm}
\setlength{\parskip}{0pt}
\setlength{\LTpre}{0pt}
\setlength{\LTpost}{0pt}
\DeclareCaptionLabelSeparator{gost}{~–~}\captionsetup{labelsep=gost, justification=centering, singlelinecheck=false}
\renewcommand{\arraystretch}{1.22}
\setlength{\tabcolsep}{4.2pt}
\setlength{\extrarowheight}{2.2pt}
\newcolumntype{L}[1]{>{\RaggedRight\arraybackslash}m{#1}}
\newcolumntype{C}[1]{>{\centering\arraybackslash}m{#1}}
\newcolumntype{R}[1]{>{\RaggedLeft\arraybackslash}m{#1}}
\setlength{\emergencystretch}{4em}
\setcounter{secnumdepth}{3}
\setcounter{tocdepth}{3}

\titleformat{\section}{\bfseries\Large}{\thesection}{0.9em}{}
\titleformat{\subsection}{\bfseries\large}{\thesubsection}{0.6em}{}
\titleformat{\subsubsection}{\bfseries\normalsize}{\thesubsubsection}{0.6em}{}

\titlespacing*{\section}{0pt}{0.9\baselineskip}{0.7\baselineskip}
\titlespacing*{\subsection}{0pt}{0.85\baselineskip}{0.55\baselineskip}
\titlespacing*{\subsubsection}{0pt}{0.7\baselineskip}{0.4\baselineskip}
\makeatletter
\renewcommand{\tableofcontents}{%
  \section*{\contentsname}%
  \@starttoc{toc}%
}
\makeatother
\renewcommand{\cftsecfont}{\normalfont}
\renewcommand{\cftsecpagefont}{\normalfont}
\newcommand{\pycmd}[1]{%
  \begin{tcolorbox}[enhanced,breakable,colback=white,colframe=black,boxrule=0.5pt,left=5pt,right=5pt,top=5pt,bottom=5pt,arc=0pt]
  {\ttfamily\small #1}
  \end{tcolorbox}%
}
\newenvironment{tighttable}[1][\small]{#1\setlength{\tabcolsep}{3.6pt}\renewcommand{\arraystretch}{1.14}}{\par\normalsize}

\newcommand{\insertfig}[2][0.86\textwidth]{%
    \IfFileExists{#2}{\includegraphics[width=#1]{#2}}{%
        \fbox{\parbox[c][6cm][c]{0.8\textwidth}{\centering Файл изображения не найден:\\ \texttt{\detokenize{#2}}}}%
    }%
}

\begin{document}


\tableofcontents
\newpage

\phantomsection
\section*{Введение}
\addcontentsline{toc}{section}{Введение}

Во второй практической работе использована та же предметная база, что и в предыдущем отчёте: официальные таблицы UN Tourism по международным туристским прибытиям и поступлениям от въездного туризма. Это позволяет продолжить анализ на той же статистической основе: в первой работе показатель был рассчитан и описан, а здесь на тех же данных проводится проверка статистических гипотез.

В центре внимания находится показатель \textit{Receipts per Arrival}, то есть средний объём поступлений на одно международное прибытие. Он позволяет увидеть туризм не только как поток поездок, но и как экономический результат. Для стран с близким числом прибытий этот показатель особенно полезен: по нему хорошо заметно, насколько различается доходность одной поездки.

Цель работы состоит в последовательном применении четырех статистических критериев к данным по странам Европы, корректной записи гипотез, интерпретации результатов и формулировке выводов. Логика исследования выстроена по шагам: сначала оценивается форма распределения, затем проверяется связь исходных показателей, после этого сравниваются независимые группы стран, а в завершение выполняется парное сопоставление за 2019 и 2021 годы.

\section{Основная часть}

\subsection{Выбор статистических критериев и дополнительных данных}

Для исследования выбраны четыре критерия разных типов:
\begin{enumerate}[leftmargin=1.25cm]
    \item Критерий Шапиро-Уилка для проверки гипотезы о нормальности распределения показателя \textit{Receipts per Arrival} за 2021 год.
    \item Ранговый коэффициент корреляции Спирмена для проверки связи между числом международных прибытий и суммарными поступлениями от въездного туризма.
    \item Критерий Манна-Уитни для сравнения показателя \textit{Receipts per Arrival} в двух независимых группах стран с очень малым и очень большим туристским потоком.
    \item Критерий Уилкоксона для связанных выборок для парного сравнения показателя \textit{Receipts per Arrival} в 2019 и 2021 годах по одним и тем же странам.
\end{enumerate}

Выбранные критерии решают четыре разные статистические задачи: проверку распределения, анализ связи между показателями, сравнение независимых групп и сравнение связанных наблюдений. Поэтому каждый следующий критерий используется для своей задачи и не дублирует предыдущий.

Помимо основной выборки за 2021 год, были подготовлены дополнительные данные:
\begin{enumerate}[leftmargin=1.25cm]
    \item Значения числа международных прибытий и поступлений за 2021 год для проверки корреляции.
    \item Две контрастные группы стран, сформированные по квартилям распределения прибытий за 2021 год.
    \item Парная выборка из 31 страны, для которых показатель \textit{Receipts per Arrival} удалось рассчитать и за 2019, и за 2021 год.
\end{enumerate}

При этом исследование не рассматривается как анализ временного ряда. Единицами наблюдения остаются страны, а сравнение 2019 и 2021 годов используется только как парное сопоставление двух срезов одной и той же совокупности.

\subsection{Описание исходных данных и их источников}

В работе использованы данные официальной базы UN Tourism из раздела \textit{Tourism Statistics Database - Inbound Tourism}. Для расчётов взяты две таблицы:
\begin{enumerate}[leftmargin=1.25cm]
    \item \textit{Inbound Tourism: Arrivals}.
    \item \textit{Inbound Tourism: Expenditure}.
\end{enumerate}

Из первой таблицы взят показатель \textit{inbound / trips / total / total / overnight visitors (tourists)}.\linebreak Он отражает количество международных прибытий туристов с ночевкой и измеряется в тысячах поездок.

Из второй таблицы взят показатель \textit{inbound / expenditure / balance of payments / total / visitors}.\linebreak Он отражает суммарные поступления от въездного туризма и измеряется в миллионах долларов США.

На их основе рассчитывается исследуемый показатель:
\begin{equation}
R_i = \frac{E_i}{A_i} \cdot 1000,
\end{equation}
где $R_i$ обозначает \textit{Receipts per Arrival} для $i$-й страны, USD на одно прибытие; $E_i$ обозначает поступления от въездного туризма, млн USD; $A_i$ обозначает международные прибытия, тыс. поездок.

Множитель $1000$ используется из-за различия исходных единиц измерения: поступления заданы в миллионах долларов, а прибытия выражены в тысячах поездок. После пересчета показатель получается в долларах США на одно прибытие.

Основные источники данных:
\begin{enumerate}[leftmargin=1.25cm]
    \item Официальный раздел UN Tourism по въездному туризму: \url{https://www.untourism.int/tourism-statistics/tourism-data-inbound-tourism}.
    \item Таблица \textit{UN\_Tourism\_inbound\_arrivals\_12\_2025.xlsx}: \url{https://pre-webunwto.s3.amazonaws.com/s3fs-public/2025-12/UN_Tourism_inbound_arrivals_12_2025.xlsx}.
    \item Таблица \textit{UN\_Tourism\_inbound\_expenditure\_12\_2025.xlsx}: \url{https://pre-webunwto.s3.amazonaws.com/s3fs-public/2025-12/UN_Tourism_inbound_expenditure_12_2025.xlsx}.
\end{enumerate}

В основную выборку за 2021 год вошла 41 страна Европы, по которым одновременно присутствуют оба показателя. Для парного сравнения 2019 и 2021 годов после отбора осталась сопоставимая выборка объемом 31 страна.

\subsection{Формальное представление подготовленных данных}

Основная подготовленная выборка приведена в таблице~\ref{tab:maindata}. В ней по каждой стране показаны число прибытий, поступления и рассчитанный показатель \textit{Receipts per Arrival}. Таблица дана полностью, потому что объем выборки еще остается удобным для чтения, а скрывать часть наблюдений здесь нецелесообразно: именно по этим странам дальше выполняются все статистические проверки.

\begin{tighttable}[\footnotesize]
\begin{longtable}{|C{0.65cm}|L{3.6cm}|R{1.95cm}|R{2.2cm}|R{2.2cm}|}
\caption{Подготовленная выборка по странам Европы за 2021 год}\label{tab:maindata}\\
\hline
№ & Страна & \makecell[r]{Прибытия,\\ тыс.} & \makecell[r]{Поступления,\\ млн USD} & \makecell[r]{Receipts per\\ Arrival, USD} \\ \hline
\endfirsthead
\hline
№ & Страна & \makecell[r]{Прибытия,\\ тыс.} & \makecell[r]{Поступления,\\ млн USD} & \makecell[r]{Receipts per\\ Arrival, USD} \\ \hline
\endhead
\endfoot
\hline
\endlastfoot
1 & Австрия & 12728,00 & 11944,47 & 938,44 \\ \hline
2 & Албания & 5515,00 & 2479,95 & 449,67 \\ \hline
3 & Андорра & 1949,00 & 1924,88 & 987,62 \\ \hline
4 & Армения & 876,00 & 813,43 & 928,57 \\ \hline
5 & Беларусь & 787,00 & 661,53 & 840,57 \\ \hline
6 & Бельгия & 3243,00 & 8305,29 & 2560,99 \\ \hline
7 & Болгария & 2300,00 & 2720,34 & 1182,76 \\ \hline
8 & Босния и Герцеговина & 502,00 & 1061,22 & 2113,98 \\ \hline
9 & Великобритания & 6287,00 & 37790,49 & 6010,89 \\ \hline
10 & Венгрия & 7929,10 & 5692,55 & 717,93 \\ \hline
11 & Германия & 11652,00 & 22137,26 & 1899,87 \\ \hline
12 & Греция & 14705,00 & 13666,11 & 929,35 \\ \hline
13 & Грузия & 1577,00 & 1350,61 & 856,44 \\ \hline
14 & Дания & 7555,00 & 4426,76 & 585,94 \\ \hline
15 & Исландия & 698,00 & 1302,80 & 1866,48 \\ \hline
16 & Испания & 31181,00 & 34176,16 & 1096,06 \\ \hline
17 & Италия & 26888,00 & 25354,93 & 942,98 \\ \hline
18 & Кипр & 1936,93 & 2099,71 & 1084,04 \\ \hline
19 & Латвия & 478,00 & 739,31 & 1546,68 \\ \hline
20 & Литва & 947,60 & 585,15 & 617,51 \\ \hline
21 & Люксембург & 615,00 & 6069,44 & 9869,00 \\ \hline
22 & Мальта & 968,00 & 3195,54 & 3301,18 \\ \hline
23 & Молдова & 68,85 & 484,07 & 7030,59 \\ \hline
24 & Нидерланды & 6248,00 & 13677,21 & 2189,05 \\ \hline
25 & Норвегия & 1435,00 & 2316,74 & 1614,45 \\ \hline
26 & Польша & 9722,00 & 10150,00 & 1044,02 \\ \hline
27 & Португалия & 9617,00 & 13904,15 & 1445,79 \\ \hline
28 & Румыния & 6788,90 & 3524,78 & 519,20 \\ \hline
29 & Сан-Марино & 94,00 & 224,76 & 2391,06 \\ \hline
30 & Северная Македония & 294,00 & 386,68 & 1315,23 \\ \hline
31 & Сербия & 871,00 & 2179,16 & 2501,90 \\ \hline
32 & Словения & 1832,00 & 2041,47 & 1114,34 \\ \hline
33 & Турция & 29925,00 & 35209,00 & 1176,57 \\ \hline
34 & Украина & 3973,00 & 1387,00 & 349,11 \\ \hline
35 & Франция & 48395,00 & 44430,27 & 918,08 \\ \hline
36 & Хорватия & 10641,00 & 10899,96 & 1024,34 \\ \hline
37 & Черногория & 1554,00 & 902,90 & 581,02 \\ \hline
38 & Чехия & 3768,00 & 3932,47 & 1043,65 \\ \hline
39 & Швейцария & 4390,00 & 13069,16 & 2977,03 \\ \hline
40 & Швеция & 2990,00 & 6056,00 & 2025,42 \\ \hline
41 & Эстония & 808,00 & 896,39 & 1109,40 \\
\end{longtable}
\end{tighttable}

Даже без формальных расчётов видно, что выборка неоднородна. В нижней части распределения есть страны, где показатель остаётся ниже 600 USD на одно прибытие, а в верхней части встречаются значения в несколько тысяч долларов. Уже по таблице заметно, что распределение асимметрично, а значит, для дальнейшего анализа разумно опираться прежде всего на непараметрические процедуры.

\subsection{Наглядное представление данных}

Сначала была построена гистограмма показателя \textit{Receipts per Arrival} (рис.~\ref{fig:hist}). Она нужна не только как иллюстрация, но и как первый статистический ориентир: по ней видно, где сосредоточена основная масса наблюдений и насколько сильно распределение тянется вправо.

\begin{figure}[H]
    \centering
    \insertfig{assets/hist_rpa_2021.png}
    \caption{Гистограмма распределения показателя \textit{Receipts per Arrival} за 2021 год}
    \label{fig:hist}
\end{figure}

По рисунку видно, что большая часть стран сосредоточена в зоне относительно умеренных значений, а несколько наблюдений сильно уходят вправо. Это визуально поддерживает гипотезу о ненормальности сырых данных и объясняет, почему далее используются ранговые критерии.

Для исходных показателей дополнительно построен график зависимости прибытия и поступления в логарифмических шкалах (рис.~\ref{fig:scatter}). Логарифмический масштаб выбран для того, чтобы на одном рисунке были хорошо видны и крупные туристские рынки, и страны с небольшими значениями.

\begin{figure}[H]
    \centering
    \insertfig{assets/scatter_arrivals_expenditure.png}
    \caption{Связь международных прибытий и поступлений от въездного туризма в 2021 году}
    \label{fig:scatter}
\end{figure}

Точки образуют восходящее облако. Уже по рисунку видно, что страны с большим туристским потоком, как правило, получают и большие совокупные поступления. Формально эта связь будет проверена критерием Спирмена.

\subsection{Описание применяемых критериев}

\subsubsection{Критерий Шапиро-Уилка}

Критерий Шапиро-Уилка используется для проверки гипотезы о нормальности распределения. Он особенно удобен на выборках небольшого и среднего объема и считается одним из наиболее чувствительных тестов на отклонение от нормального закона.

Статистика критерия имеет вид
\begin{equation}
W = \frac{\left(\sum\limits_{i=1}^{n} a_i x_{(i)}\right)^2}{\sum\limits_{i=1}^{n} (x_i-\overline{x})^2},
\end{equation}
где $x_{(i)}$ обозначают упорядоченные наблюдения, а $a_i$ обозначают коэффициенты, зависящие от нормального распределения и объема выборки.

В данной работе критерий применяется к выборке значений \textit{Receipts per Arrival} за 2021 год ($n=41$).

\par\indent Нулевая гипотеза:
\begin{equation*}
H_0: \text{распределение } \textit{Receipts per Arrival} \text{ не отличается от нормального.}
\end{equation*}
\par\indent Альтернативная гипотеза:
\begin{equation*}
H_1: \text{распределение } \textit{Receipts per Arrival} \text{ не является нормальным.}
\end{equation*}

Команда в Python:
\pycmd{stats.shapiro(rpa\_2021)}
Здесь \texttt{rpa\_2021} обозначает одномерный массив значений показателя за 2021 год.

\subsubsection{Критерий Спирмена}

Коэффициент Спирмена применяется для проверки монотонной связи между двумя количественными признаками, когда нет оснований опираться на нормальность распределений и линейность связи.

Для случая без совпадающих рангов коэффициент вычисляется по формуле
\begin{equation}
\rho_s = 1 - \frac{6\sum\limits_{i=1}^{n} d_i^2}{n(n^2-1)},
\end{equation}
где $d_i$ обозначает разность рангов двух признаков у $i$-го наблюдения.

В работе сравниваются два исходных показателя за 2021 год: число международных прибытий и суммарные поступления от въездного туризма ($n=41$). Это позволяет проверить, насколько согласованно изменяются базовые туристские характеристики стран.

\par\indent Нулевая гипотеза:
\begin{equation*}
H_0: \rho_s = 0, \text{ монотонная связь между прибытиями и поступлениями отсутствует.}
\end{equation*}
\par\indent Альтернативная гипотеза:
\begin{equation*}
H_1: \rho_s \neq 0, \text{ монотонная связь присутствует.}
\end{equation*}

Команда в Python:
\pycmd{stats.spearmanr(arrivals\_2021, expenditure\_2021, alternative='two-sided')}
Параметр \texttt{alternative='two-sided'} означает двустороннюю проверку гипотезы.

\subsubsection{Критерий Манна-Уитни}

Критерий Манна-Уитни предназначен для сравнения двух независимых выборок, когда данные можно упорядочить по величине, но нет оснований предполагать нормальность распределений. Он фактически сравнивает ранговую структуру двух групп.

Одна из форм записи статистики имеет вид
\begin{equation}
U_1 = n_1n_2 + \frac{n_1(n_1+1)}{2} - R_1,
\end{equation}
где $R_1$ обозначает сумму рангов первой группы. Аналогично вычисляется $U_2$, а в качестве статистики берут меньшее из двух значений.

В данной работе формируются две контрастные независимые группы: страны из нижнего квартиля по числу прибытий туристов ($n_1=11$) и страны из верхнего квартиля ($n_2=11$). Такой прием выбран специально: он позволяет сравнить не соседние по величине наблюдения, а действительно разные по масштабу туристского потока страны.

\par\indent Нулевая гипотеза:
\begin{equation*}
H_0: \text{распределения } \textit{Receipts per Arrival} \text{ в двух группах одинаковы.}
\end{equation*}
\par\indent Альтернативная гипотеза:
\begin{equation*}
H_1: \text{распределения } \textit{Receipts per Arrival} \text{ в двух группах различаются.}
\end{equation*}

Команда в Python:
\pycmd{stats.mannwhitneyu(rpa\_low, rpa\_high, alternative='two-sided', method='asymptotic')}
Здесь \texttt{rpa\_low} и \texttt{rpa\_high} обозначают массивы значений показателя в крайних группах, а параметр \texttt{method='asymptotic'} задает асимптотический расчет $p$-значения.

\subsubsection{Критерий Уилкоксона для связанных выборок}

Критерий Уилкоксона используется для сравнения двух связанных выборок. Он основан на ранжировании модулей попарных различий и учитывает направление сдвига.

Пусть $d_i = x_i - y_i$ обозначает разность внутри пары. После исключения нулевых различий ранжируются значения $|d_i|$, а статистика строится по сумме рангов различий одного знака:
\begin{equation}
W = \sum \text{rank}(|d_i|) \quad \text{для } d_i > 0.
\end{equation}

В работе сравниваются значения \textit{Receipts per Arrival} за 2019 и 2021 годы для одних и тех же 31 стран. Это связанная выборка, поэтому обычное сравнение независимых групп здесь было бы некорректным.

\par\indent Нулевая гипотеза:
\begin{equation*}
H_0: \text{медиана попарных различий } (R_{2021}-R_{2019}) \text{ равна нулю.}
\end{equation*}
\par\indent Альтернативная гипотеза:
\begin{equation*}
H_1: \text{медиана попарных различий } (R_{2021}-R_{2019}) > 0.
\end{equation*}

Команда в Python:
\pycmd{stats.wilcoxon(\\
\quad rpa\_2021\_paired,\\
\quad rpa\_2019\_paired,\\
\quad alternative='greater',\\
\quad zero\_method='wilcox',\\
\quad correction=False,\\
\quad method='auto'\\
)}
Параметр \texttt{alternative='greater'} соответствует проверке именно на рост показателя в 2021 году. Параметр \texttt{zero\_method='wilcox'} исключает нулевые разности, а \texttt{method='auto'} позволяет автоматически выбрать способ вычисления $p$-значения.

\subsection{Результаты проверки критериев и их интерпретация}

Итоговые результаты сведены в таблицу~\ref{tab:criteria}.

\begin{table}[H]
\centering
\begin{tighttable}[\footnotesize]
\caption{Сводные результаты статистических критериев}
\label{tab:criteria}
\begin{tabular}{|L{3.0cm}|C{1.95cm}|C{2.3cm}|C{1.55cm}|L{4.55cm}|}
\hline
Критерий & Объем выборки & Статистика & $p$-value & Вывод \\ \hline
Шапиро-Уилка & 41 & $W=0,627$ & $< 0{,}001$ & распределение не является нормальным \\ \hline
Спирмена & 41 & $\rho_s=0,871$ & $< 0{,}001$ & есть значимая положительная связь \\ \hline
Манна-Уитни & 11 и 11 & $U=95,0$ & 0,026 & группы статистически различаются \\ \hline
Уилкоксона & 31 пар & $W=496,0$ & $< 0{,}001$ & показатель в 2021 году выше \\ \hline
\end{tabular}
\end{tighttable}
\end{table}

\subsubsection{Критерий Шапиро-Уилка}
Для основной выборки за 2021 год получено $W=0,627$ и $p<0{,}001$. На уровне значимости $\alpha=0,05$ нулевая гипотеза отвергается. Следовательно, распределение сырых значений \textit{Receipts per Arrival} нельзя считать нормальным.

Этот результат хорошо согласуется и с гистограммой: справа есть длинный хвост, а несколько стран заметно выбиваются из основной массы наблюдений. Для логарифмов показателя дополнительная проверка даёт $p=0{,}055$, то есть логарифмирование заметно сглаживает асимметрию. Основной вывод от этого не меняется, но сама форма распределения становится понятнее.

\subsubsection{Критерий Спирмена}
Для связи между международными прибытиями и суммарными поступлениями получено $\rho_s=0,871$ и $p<0{,}001$. Нулевая гипотеза об отсутствии монотонной связи отвергается. Значит, между двумя исходными туристскими показателями существует сильная положительная связь.

Этот результат вполне ожидаем: страны, которые принимают больше туристов, обычно получают и большие совокупные туристские поступления. Критерий Спирмена подтверждает эту связь статистически.

\subsubsection{Критерий Манна-Уитни}
Нижний квартиль по числу прибытий туристов составил 876,00 тыс. поездок, а верхний квартиль составил 7929,10 тыс. поездок. После отбора были получены две независимые группы по 11 стран в каждой. Для них рассчитано $U=95,0$ и $p=0{,}026$. На уровне значимости $0,05$ нулевая гипотеза отвергается.

Это означает, что распределение \textit{Receipts per Arrival} у стран с очень малым и очень большим потоком туристов различается статистически значимо. Медиана показателя в группе стран с малым потоком равна 1866,48 USD, а в группе стран с большим потоком равна 1024,34 USD. Иначе говоря, для группы стран с небольшим туристским потоком характерна более высокая медиана показателя, чем для крупных массовых направлений.

Наглядно это видно на диаграмме размаха (рис.~\ref{fig:groups}).

\begin{figure}[H]
    \centering
    \insertfig{assets/boxplot_flow_groups.png}
    \caption{Сравнение \textit{Receipts per Arrival} в крайних группах стран по числу прибытий}
    \label{fig:groups}
\end{figure}

\subsubsection{Критерий Уилкоксона}
В парной выборке объемом 31 страна получено $W=496,0$ и $p<0{,}001$ при односторонней альтернативе о росте показателя в 2021 году. Нулевая гипотеза отвергается. Следовательно, показатель \textit{Receipts per Arrival} в 2021 году статистически значимо выше уровня 2019 года.

Этот результат хорошо подтверждается составом парной выборки: среди рассмотренных стран не оказалось ни одной, где рассчитанный показатель за 2021 год был бы ниже значения 2019 года. Минимальный положительный прирост составил 22,89 USD на одно прибытие, а максимальный прирост составил 4238,17 USD.

По сводным характеристикам тот же сдвиг виден достаточно ясно: среднее значение по сопоставимой выборке выросло с 1042,47 до 1733,44 USD, а медиана выросла с 778,60 до 1084,04 USD. На рисунке~\ref{fig:paired} этот результат показан графически.

\begin{figure}[H]
    \centering
    \insertfig[0.78\textwidth]{assets/paired_boxplot_2019_2021.png}
    \caption{Сравнение распределения \textit{Receipts per Arrival} в 2019 и 2021 годах по сопоставимой выборке стран}
    \label{fig:paired}
\end{figure}

\subsection{Общие выводы по основной части}

Проверка четырех критериев показала, что результаты основной части хорошо согласуются между собой. Сначала было установлено, что сырые значения \textit{Receipts per Arrival} не подчиняются нормальному закону. Затем была подтверждена сильная положительная связь между международными прибытиями и суммарными поступлениями. После этого выяснилось, что страны с очень малым и очень большим туристским потоком различаются по величине поступлений на одно прибытие, а в конце парное сравнение показало статистически значимый рост показателя в 2021 году по отношению к 2019 году.

Это означает, что туристский рынок Европы в 2021 году был неоднородным; крупные рынки действительно собирали больше денег в абсолютном выражении, но не обязательно лидировали по поступлениям на одного туриста; кроме того, после пандемийного шока доходность одной поездки в сопоставимой выборке стран оказалась выше, чем в 2019 году.


\section{Расширение исследования}

\subsection{Логика последовательного применения критериев}

Расширенная часть построена как продолжение основной части, а не как набор отдельных проверок. Последовательность применения критериев приведена в таблице~\ref{tab:logic}.

\begin{table}[H]
\centering
\begin{tighttable}[\footnotesize]
\caption{Логика последовательного применения критериев}
\label{tab:logic}
\begin{tabular}{|C{0.8cm}|L{2.8cm}|L{3.75cm}|L{4.2cm}|}
\hline
Этап & Критерий & Полученный результат & Методическое следствие \\ \hline
1 & Шапиро-Уилка & $W=0{,}627$, $p<0{,}001$; нормальность сырых значений не подтверждается & Для дальнейших проверок целесообразно опираться на ранговые непараметрические критерии \\ \hline
2 & Спирмена & $\rho_s=0{,}871$, $p<0{,}001$; между прибытиями и поступлениями есть сильная положительная связь & После проверки формы распределения исследование переходит к анализу связи исходных переменных \\ \hline
3 & Манна-Уитни & $U=95{,}0$, $p=0{,}026$; крайние группы стран статистически различаются & После анализа связи можно отдельно сравнить независимые группы, различающиеся по масштабу туристского потока \\ \hline
4 & Уилкоксона & $W=496{,}0$, $p<0{,}001$; показатель в 2021 году выше, чем в 2019 году & Заключительный этап посвящён уже не группам стран, а парному межгодовому сдвигу по одной и той же совокупности \\ \hline
\end{tabular}
\end{tighttable}
\end{table}

Такая последовательность выглядит логичной. Сначала проверяется, можно ли считать сырые данные нормальными. После отклонения этой гипотезы становится ясно, что дальше лучше использовать непараметрические процедуры. Затем отдельно рассматриваются связь исходных переменных и различия между независимыми группами. После этого выполняется парное сопоставление 2019 и 2021 годов, где сравниваются одни и те же страны. Таким образом, каждый критерий отвечает на свой вопрос.

\subsection{Дополнительное применение критериев к уже выбранным данным}

Для проверки устойчивости выводов к уже выбранным наборам данных были применены дополнительные критерии. Их результаты сведены в таблицу~\ref{tab:extra_criteria}.

\begin{table}[H]
\centering
\begin{tighttable}[\footnotesize]
\caption{Дополнительные критерии, применённые к уже выбранным наборам данных}
\label{tab:extra_criteria}
\begin{tabular}{|L{3.1cm}|L{3.6cm}|C{1.9cm}|C{1.55cm}|L{3.6cm}|}
\hline
Критерий & Набор данных & Статистика & $p$-value & Интерпретация \\ \hline
Шапиро-Уилка для $\ln R_i$ & Логарифмы \textit{Receipts per Arrival}, 2021 ($n=41$) & $W=0{,}947$ & 0,055 & После логарифмирования отклонение от нормальности на уровне 0,05 уже не фиксируется \\ \hline
Кендалла & Прибытия и поступления, 2021 ($n=41$) & $\tau=0{,}715$ & $< 0{,}001$ & Сильная положительная монотонная связь подтверждается и вторым ранговым коэффициентом \\ \hline
Точный критерий знаков & Пары 2019/2021 ($n=31$) & $S=31$ & $< 0{,}001$ & Во всех 31 парах значение 2021 года оказалось выше значения 2019 года \\ \hline
\end{tabular}
\end{tighttable}
\end{table}

Дополнительная проверка логарифмов показывает, что основная проблема исходного распределения связана с правым хвостом и несколькими очень крупными наблюдениями. После перехода к $\ln R_i$ распределение становится заметно более ровным: значение $p=0{,}055$ уже не даёт оснований отвергать нормальность на уровне значимости 0,05. Следовательно, асимметрия исходного распределения обусловлена прежде всего несколькими крупными значениями в правой части распределения.

Критерий Кендалла применён к тем же данным, что и коэффициент Спирмена. Полученное значение $\tau=0{,}715$ при $p<0{,}001$ повторяет основной вывод о сильной положительной монотонной связи между прибытием туристов и суммарными поступлениями. Это означает, что вывод о наличии связи не зависит от выбора конкретного рангового коэффициента.

Точный критерий знаков дополнил критерий Уилкоксона для парной выборки. Если критерий Уилкоксона учитывает и направление, и величину попарных сдвигов, то критерий знаков фиксирует только направление изменения. В рассматриваемой выборке все 31 разность положительны, поэтому дополнительная проверка полностью подтверждает рост \textit{Receipts per Arrival} в 2021 году по сравнению с 2019 годом.

\subsection{Дополнительные данные и построение выборок}

Для расширенной части были подготовлены дополнительные наборы данных и уточнены правила их формирования. Сводка приведена в таблице~\ref{tab:extra_data}.

\begin{table}[H]
\centering
\begin{tighttable}[\footnotesize]
\caption{Дополнительные наборы данных, использованные в расширенной части}
\label{tab:extra_data}
\begin{tabular}{|L{3.4cm}|L{3.15cm}|C{1.5cm}|L{3.2cm}|}
\hline
Набор данных & Правило формирования & Объём & Назначение \\ \hline
Логарифмы $\ln R_i$ за 2021 год & Логарифмирование всех 41 исходных значений \textit{Receipts per Arrival} & 41 страна & Уточнение формы распределения и проверка влияния правого хвоста \\ \hline
Данные по прибытиям и поступлениям за 2021 год & Совмещение двух исходных таблиц UN Tourism по одинаковому набору стран & 41 страна & Проверка силы и направления связи между базовыми туристскими показателями \\ \hline
Контрастные группы стран по числу прибытий & Отбор стран не выше $Q_1=876{,}00$ тыс. и не ниже $Q_3=7929{,}10$ тыс. прибытий & 11 и 11 стран & Сравнение независимых групп стран с очень малым и очень большим туристским потоком \\ \hline
Парная выборка 2019/2021 & Отбор только тех стран, для которых показатель рассчитан в обоих годах & 31 страна & Корректное сопоставление межгодового сдвига по одной и той же совокупности \\ \hline
\end{tabular}
\end{tighttable}
\end{table}

Таким образом, из исходных данных были сформированы несколько выборок, каждая для своей статистической задачи. Это позволило отдельно провести сравнение независимых групп, корреляционный анализ и парное сопоставление.

В таблице~\ref{tab:groups} приведён состав контрастных групп стран по числу прибытий, использованных в критерии Манна-Уитни. Для компактности в заголовке таблицы используется сокращение RPA, соответствующее показателю \textit{Receipts per Arrival}.

\begin{tighttable}[\footnotesize]
\begin{longtable}{|C{0.6cm}|L{2.45cm}|R{1.42cm}|R{1.55cm}|L{2.45cm}|R{1.42cm}|R{1.55cm}|}
\caption{Состав крайних групп стран по числу прибытий за 2021 год}\label{tab:groups}\\
\hline
\multirow{2}{*}{\centering №} &
\multicolumn{3}{c|}{\makecell[c]{Нижний квартиль\\по числу прибытий}} &
\multicolumn{3}{c|}{\makecell[c]{Верхний квартиль\\по числу прибытий}} \\ \cline{2-7}
& \multicolumn{1}{c|}{Страна}
& \multicolumn{1}{c|}{\makecell[c]{Прибытия,\\ тыс.}}
& \multicolumn{1}{c|}{\makecell[c]{RPA,\\ USD}}
& \multicolumn{1}{c|}{Страна}
& \multicolumn{1}{c|}{\makecell[c]{Прибытия,\\ тыс.}}
& \multicolumn{1}{c|}{\makecell[c]{RPA,\\ USD}} \\ \hline
\endfirsthead

\hline
\multirow{2}{*}{\centering №} &
\multicolumn{3}{c|}{\makecell[c]{Нижний квартиль\\по числу прибытий}} &
\multicolumn{3}{c|}{\makecell[c]{Верхний квартиль\\по числу прибытий}} \\ \cline{2-7}
& \multicolumn{1}{c|}{Страна}
& \multicolumn{1}{c|}{\makecell[c]{Прибытия,\\ тыс.}}
& \multicolumn{1}{c|}{\makecell[c]{RPA,\\ USD}}
& \multicolumn{1}{c|}{Страна}
& \multicolumn{1}{c|}{\makecell[c]{Прибытия,\\ тыс.}}
& \multicolumn{1}{c|}{\makecell[c]{RPA,\\ USD}} \\ \hline
\endhead

\hline
\endfoot

\hline
\endlastfoot

1 & Люксембург & 615,00 & 9869,00 & Германия & 11652,00 & 1899,87 \\ \hline
2 & Молдова & 68,85 & 7030,59 & Португалия & 9617,00 & 1445,79 \\ \hline
3 & Сербия & 871,00 & 2501,90 & Турция & 29925,00 & 1176,57 \\ \hline
4 & Сан-Марино & 94,00 & 2391,06 & Испания & 31181,00 & 1096,06 \\ \hline
5 & Босния и Герцеговина & 502,00 & 2113,98 & Польша & 9722,00 & 1044,02 \\ \hline
6 & Исландия & 698,00 & 1866,48 & Хорватия & 10641,00 & 1024,34 \\ \hline
7 & Латвия & 478,00 & 1546,68 & Италия & 26888,00 & 942,98 \\ \hline
8 & Северная Македония & 294,00 & 1315,23 & Австрия & 12728,00 & 938,44 \\ \hline
9 & Эстония & 808,00 & 1109,40 & Греция & 14705,00 & 929,35 \\ \hline
10 & Армения & 876,00 & 928,57 & Франция & 48395,00 & 918,08 \\ \hline
11 & Беларусь & 787,00 & 840,57 & Венгрия & 7929,10 & 717,93 \\
\end{longtable}
\end{tighttable}

\subsection{Подробные выводы по расширенной части}

По результатам расширенной части можно сделать следующие выводы.

\begin{enumerate}[leftmargin=1.25cm]
    \item Дополнительная проверка логарифмов показала, что основное отклонение исходного распределения от нормальности связано с длинным правым хвостом. Для сырых значений гипотеза о нормальности уверенно отвергается, тогда как для $\ln R_i$ получено $p=0{,}055$. Значит, после логарифмирования распределение становится заметно ближе к нормальному, а выраженная асимметрия исходной выборки в первую очередь объясняется несколькими очень крупными значениями.
    \item Критерий Кендалла подтвердил тот же вывод, что и коэффициент Спирмена: между числом международных прибытий и суммарными поступлениями существует сильная положительная монотонная связь. Следовательно, рост совокупных поступлений по странам Европы действительно сопровождается ростом туристского потока, и этот вывод не зависит от выбора конкретного рангового коэффициента.
    \item Контрастные группы и парная выборка понадобились для решения разных статистических задач. Критерий Манна-Уитни сравнивает две независимые группы стран, различающиеся по масштабу туристского потока в одном и том же году, тогда как парная выборка из 31 страны позволяет сопоставить 2019 и 2021 годы без смены состава наблюдений. Благодаря этому межгрупповое сравнение и межгодовой сдвиг не смешиваются в одном расчёте.
    \item Расширенная часть подтверждает, что европейский туристский рынок в 2021 году был неоднородным сразу по нескольким признакам. Страны различались по масштабу потока, по абсолютной сумме поступлений и по уровню поступлений на одно прибытие. Кроме того, сопоставление с 2019 годом показывает, что восстановление происходило неравномерно: в сопоставимой выборке показатель \textit{Receipts per Arrival} вырос у всех стран, то есть в 2021 году на одно международное прибытие в среднем приходилось больше поступлений, чем до кризисного периода.
\end{enumerate}

Таким образом, дополнительные проверки подтвердили основные выводы и уточнили их. Они показали, с чем связана асимметрия распределения, подтвердили вывод о корреляции и отдельно показали различия между группами стран и межгодовое изменение показателя.

\clearpage
\phantomsection
\section*{Заключение}
\addcontentsline{toc}{section}{Заключение}

В ходе практической работы была выполнена статистическая проверка четырех критериев на данных UN Tourism по странам Европы. Исследование опиралось на тот же предметный материал, что и предыдущая практическая работа, но теперь данные использовались уже не только для описания, а для проверки гипотез.

В результате установлено следующее:
\begin{enumerate}[leftmargin=1.25cm]
    \item Распределение показателя \textit{Receipts per Arrival} за 2021 год не является нормальным.
    \item Между международными прибытиями туристов и суммарными поступлениями существует сильная положительная монотонная связь.
    \item Страны с очень малым и очень большим туристским потоком статистически значимо различаются по уровню поступлений на одно прибытие.
    \item В сопоставимой выборке из 31 страны показатель \textit{Receipts per Arrival} в 2021 году статистически значимо выше, чем в 2019 году.
\end{enumerate}

Таким образом, выбранные критерии дополнили друг друга и позволили рассмотреть данные с разных сторон. Полученные результаты показывают, что статистические методы нужны не только для проверки гипотез, но и для содержательной интерпретации исходных данных. За одинаковой величиной туристского потока могут скрываться разные модели доходности, а сравнение по годам помогает увидеть изменения в самой структуре рынка.

\clearpage
\phantomsection
\section*{Приложение А. Парная выборка 2019/2021}
\addcontentsline{toc}{section}{Приложение А. Парная выборка 2019/2021}

В таблице~\ref{tab:pairedfull} приведена полная парная выборка стран, использованная в критерии Уилкоксона.

\begin{tighttable}[\footnotesize]
\begin{longtable}{|C{0.65cm}|L{3.75cm}|R{1.9cm}|R{1.9cm}|R{1.95cm}|}
\caption{Показатель \textit{Receipts per Arrival} в 2019 и 2021 годах по сопоставимой выборке стран}\label{tab:pairedfull}\\
\hline
№ & Страна & \makecell[r]{2019,\\ USD} & \makecell[r]{2021,\\ USD} & \makecell[r]{Изменение,\\ USD} \\ \hline
\endfirsthead
\hline
№ & Страна & \makecell[r]{2019,\\ USD} & \makecell[r]{2021,\\ USD} & \makecell[r]{Изменение,\\ USD} \\ \hline
\endhead
\endfoot
\hline
\endlastfoot
1 & Австрия & 813,32 & 938,44 & 125,12 \\ \hline
2 & Албания & 401,11 & 449,67 & 48,56 \\ \hline
3 & Андорра & 618,12 & 987,62 & 369,50 \\ \hline
4 & Армения & 819,96 & 928,57 & 108,61 \\ \hline
5 & Беларусь & 584,77 & 840,57 & 255,80 \\ \hline
6 & Бельгия & 1125,55 & 2560,99 & 1435,44 \\ \hline
7 & Болгария & 621,48 & 1182,76 & 561,28 \\ \hline
8 & Босния и Герцеговина & 1022,54 & 2113,98 & 1091,44 \\ \hline
9 & Венгрия & 603,71 & 717,93 & 114,22 \\ \hline
10 & Германия & 1056,01 & 1899,87 & 843,86 \\ \hline
11 & Греция & 733,79 & 929,35 & 195,57 \\ \hline
12 & Грузия & 698,96 & 856,44 & 157,48 \\ \hline
13 & Италия & 804,64 & 942,98 & 138,34 \\ \hline
14 & Кипр & 817,45 & 1084,04 & 266,59 \\ \hline
15 & Люксембург & 5630,84 & 9869,00 & 4238,17 \\ \hline
16 & Молдова & 2995,98 & 7030,59 & 4034,61 \\ \hline
17 & Нидерланды & 1098,62 & 2189,05 & 1090,43 \\ \hline
18 & Норвегия & 1257,02 & 1614,45 & 357,43 \\ \hline
19 & Польша & 756,74 & 1044,02 & 287,29 \\ \hline
20 & Португалия & 1422,90 & 1445,79 & 22,89 \\ \hline
21 & Северная Македония & 529,02 & 1315,23 & 786,21 \\ \hline
22 & Сербия & 1082,84 & 2501,90 & 1419,07 \\ \hline
23 & Словения & 713,12 & 1114,34 & 401,21 \\ \hline
24 & Турция & 894,46 & 1176,57 & 282,12 \\ \hline
25 & Украина & 193,11 & 349,11 & 156,00 \\ \hline
26 & Франция & 778,60 & 918,08 & 139,47 \\ \hline
27 & Хорватия & 689,85 & 1024,34 & 334,48 \\ \hline
28 & Черногория & 508,76 & 581,02 & 72,25 \\ \hline
29 & Чехия & 543,79 & 1043,65 & 499,86 \\ \hline
30 & Швейцария & 1808,94 & 2977,03 & 1168,09 \\ \hline
31 & Эстония & 690,65 & 1109,40 & 418,75 \\
\end{longtable}
\end{tighttable}

\clearpage
\phantomsection
\section*{Приложение Б. Ссылка на код}
\addcontentsline{toc}{section}{Приложение Б. Ссылка на код}

\begin{enumerate}[leftmargin=1.25cm]
    \item Код анализа размещён в репозитории GitHub. URL: \url{https://github.com/Echways/SPoI}.
\end{enumerate}

\end{document}
